\chapter{Кристаллическая решётка}\label{ch:lattice}

Атомы металла, притягиваясь ненаправленно, укладываются как можно плотнее и в правильном порядке. Порядок
повторяется в пространстве --- это \textbf{кристалл}. Чтобы говорить о плоскостях залегания гидридов, о текстуре
трубы и о дифракции, нужен язык описания такого порядка: ячейки, индексы, проекции.

\section{Решётка, базис, ячейка}

\textbf{Кристаллическая решётка} --- бесконечный набор точек (узлов), получаемых из одной трансляциями
$\mathbf R = n_1\mathbf a_1 + n_2\mathbf a_2 + n_3\mathbf a_3$ с целыми $n_i$. Векторы $\mathbf a_i$ задают
\textbf{элементарную ячейку}. Кристалл --- это решётка плюс \textbf{базис}: группа атомов, «привязанная» к
каждому узлу. У простых металлов базис --- один-два атома, у соединений --- больше.

По симметрии ячеек решётки делят на 7 \textbf{сингоний} (кубическая, гексагональная, тетрагональная,
ромбоэдрическая, ромбическая, моноклинная, триклинная) и 14 \textbf{решёток Браве}. Для металлов важны четыре
структуры (рис.~\ref{fig:l1_cells}).

\fig{l1_cells}{Элементарные ячейки металлов: простая кубическая (ПК), объёмноцентрированная (ОЦК),
гранецентрированная (ГЦК) и гексагональная плотноупакованная (ГПУ). Под каждой --- число атомов на ячейку,
координационное число $z$ (число ближайших соседей) и доля объёма, занятая шарами, $f$.}

\begin{center}\small
\begin{tabular}{lcccl}
\toprule
структура & атомов на ячейку & $z$ & $f$ & связь параметра с радиусом атома $R$ \\
\midrule
ПК & 1 & 6 & 0.52 & $a = 2R$ \\
ОЦК & 2 & 8 & 0.68 & $a = 4R/\sqrt3$ (касание по диагонали куба) \\
ГЦК & 4 & 12 & 0.74 & $a = 2\sqrt2\,R$ (касание по диагонали грани) \\
ГПУ & 6 (2 на примитивную) & 12 & 0.74 & $a = 2R$, $c = \sqrt{8/3}\,a\approx1.633\,a$ (идеально) \\
\bottomrule
\end{tabular}
\end{center}

Зная ячейку, легко найти плотность: $\rho = nA/(V_{\text{яч}}N_A)$, где $n$ --- число атомов в ячейке, $A$ ---
атомная масса. Для $\alpha$-Zr объём гексагональной ячейки с двумя атомами $V = \tfrac{\sqrt3}{2}a^2c =
4.66\cdot10^{-29}$~м$^3$, откуда $\rho\approx6.5$~г/см$^3$ --- как и измеряют.

\begin{zrbox}
Цирконий \textbf{полиморфен}. Ниже 863\degC{} устойчива ГПУ фаза $\alpha$, выше --- ОЦК фаза $\beta$ (её
параметр около 0.36~нм). Под высоким давлением появляется третья, $\omega$-фаза. Превращение $\alpha\leftrightarrow\beta$
идёт с ориентационным соотношением Бюргерса $(0001)_\alpha\parallel\{110\}_\beta$,
$\langle11\bar20\rangle_\alpha\parallel\langle111\rangle_\beta$: плотнейшая плоскость ГПУ ложится на плотнейшую
плоскость ОЦК. Легирующие элементы сдвигают границу $\alpha/\beta$: кислород и олово стабилизируют $\alpha$,
ниобий, железо и водород --- $\beta$ (глава~\ref{ch:diagrams}).
\end{zrbox}

\section{Плотная упаковка: ABAB и ABCABC}

Положим на стол слой одинаковых шаров плотнейшим образом --- каждый касается шести соседей (слой A на
рис.~\ref{fig:l2_stacking}а). Над ним есть два набора ямок, B и C. Следующий слой ложится в одни из них, скажем, B.
Третий слой можно положить либо снова над A, либо в позиции C. Отсюда две плотнейшие упаковки с одинаковой
долей заполнения 0.74 и одинаковым числом соседей 12:
\begin{itemize}
\item \textbf{ABABAB\ldots} --- гексагональная плотноупакованная (ГПУ): Zr, Ti, Mg, Zn;
\item \textbf{ABCABC\ldots} --- гранецентрированная кубическая (ГЦК): Cu, Al, Ni, $\gamma$-Fe. Плотные слои --- плоскости $\{111\}$.
\end{itemize}

\fig{l2_stacking}{а) Плотный слой A и два набора ямок над ним, B и C. б) Укладка ABAB --- ГПУ (как в $\alpha$-Zr).
в) Укладка ABCABC --- ГЦК (как подрешётка циркония в $\delta$-гидриде).}

\begin{idea}
ГПУ и ГЦК отличаются не плотностью и не ближайшими соседями, а только тем, где лежит каждый третий слой. Поэтому
перейти из одной в другую можно, сдвинув каждый второй слой на треть межатомного расстояния. Ровно так подрешётка
циркония перестраивается из ГПУ $\alpha$-Zr в ГЦК $\delta$-гидрида (глава~\ref{ch:defects}): базисная плоскость
$(0001)_\alpha$ становится плоскостью $\{111\}_\delta$. Отсюда и ориентационное соотношение гидрида, и то, что
пластинка лежит около базиса.
\end{idea}

\section{Индексы Миллера}

Направление в кристалле обозначают координатами вектора в базисе ячейки, приведёнными к наименьшим целым:
$[uvw]$. Отрицательный индекс пишут с чертой: $[1\bar10]$. Все эквивалентные по симметрии направления --- семейство
$\langle uvw\rangle$.

Плоскость обозначают \textbf{индексами Миллера} $(hkl)$:
\begin{enumerate}
\item найдите, где плоскость пересекает оси, в единицах параметров ячейки (например, $1, \infty, \infty$);
\item возьмите обратные величины ($1, 0, 0$);
\item приведите к наименьшим целым: $(100)$.
\end{enumerate}
Семейство эквивалентных плоскостей --- $\{hkl\}$. В кубической решётке направление $[hkl]$ перпендикулярно
плоскости $(hkl)$, а расстояние между соседними плоскостями семейства
\begin{equation}
d_{hkl} = \frac{a}{\sqrt{h^2+k^2+l^2}} .
\end{equation}
Чем больше индексы, тем плотнее друг к другу лежат плоскости и тем меньше атомов на каждой. Плотнейшие плоскости
--- с малыми индексами: $\{111\}$ в ГЦК, $\{110\}$ в ОЦК.

\fig{l3_miller}{Плоскости $(100)$, $(110)$ и $(111)$ кубической ячейки и перпендикулярные им направления.}

\section{Гексагональные индексы Миллера--Браве}

В ГПУ удобно брать три оси $\mathbf a_1, \mathbf a_2, \mathbf a_3$ в базисной плоскости под 120\textdegree{} друг к
другу и ось $\mathbf c$ перпендикулярно. Плоскость получает четыре индекса $(hkil)$, из которых третий лишний:
$i = -(h+k)$. Зачем он? С ним плоскости, эквивалентные по симметрии, и выглядят одинаково: шесть боковых граней
призмы --- это $(10\bar10)$, $(01\bar10)$, $(\bar1100)$ и т.\,д., перестановки одних и тех же чисел. В трёхиндексной
записи они выглядели бы как $(100)$ и $(\bar110)$ --- непохоже.

Направления в четырёхиндексной записи $[uvtw]$ переводятся из трёхиндексных $[UVW]$ (по осям $\mathbf a_1$,
$\mathbf a_2$, $\mathbf c$) так:
\begin{equation}
u = \tfrac13(2U - V),\quad v = \tfrac13(2V - U),\quad t = -(u+v),\quad w = W .
\end{equation}
Так, направление вдоль $\mathbf a_1$ ($[100]$) --- это $\tfrac13[2\bar1\bar10]$; отсюда обозначение
$\langle11\bar20\rangle$ для плотнейших направлений базиса. Межплоскостное расстояние в ГПУ
\begin{equation}
\frac{1}{d^2_{hkil}} = \frac43\,\frac{h^2+hk+k^2}{a^2} + \frac{l^2}{c^2},
\end{equation}
а угол $\varphi$ плоскости с базисом (и нормали с осью $c$) даёт
$\operatorname{tg}\varphi = \frac{2}{\sqrt3}\frac{c}{a}\sqrt{h^2+hk+k^2}/l$. В отличие от кубической решётки,
направление $[hkil]$ в общем случае \emph{не} перпендикулярно плоскости $(hkil)$ --- только для базиса и призм.

\section{Межузельные пустоты}

Между плотно уложенными шарами остаются пустоты двух видов (рис.~\ref{fig:l4_sites}а).
\begin{itemize}
\item \textbf{Октаэдрическая} --- между шестью атомами (вершины октаэдра). В ГЦК их 4 на ячейку (в центре куба и
  на серединах рёбер), то есть одна на атом. Радиус вписанного шара $0.414R$.
\item \textbf{Тетраэдрическая} --- между четырьмя атомами. В ГЦК их 8 на ячейку (центры восьми малых кубиков),
  две на атом. Радиус $0.225R$.
\end{itemize}
В ГПУ то же: на атом одна октаэдрическая и две тетраэдрические пустоты. У циркония $R\approx0.16$~нм, так что
октаэдрическая пустота вмещает шар радиусом около 0.066~нм, тетраэдрическая --- 0.036~нм. Атомы внедрения больше
этих размеров и раздвигают соседей; отсюда парциальный объём водорода и его упругое взаимодействие с напряжениями
(глава~\ref{ch:diffusion}).

\fig{l4_sites}{а) Октаэдрические (жёлтые) и тетраэдрические (зелёные) пустоты ГЦК ячейки. б) Решётка
$\delta$-гидрида: ГЦК подрешётка циркония, водород в тетраэдрических пустотах (тип флюорита; при заполнении всех
пустот --- ZrH$_2$).}

\begin{zrbox}
Водород в $\alpha$-Zr занимает тетраэдрические пустоты (это установлено нейтронографией). В $\delta$-гидриде
подрешётка циркония --- ГЦК, и водород сидит в её тетраэдрических пустотах, как фтор в CaF$_2$
(рис.~\ref{fig:l4_sites}б). Тетраэдрических пустот две на атом циркония, поэтому предельный состав --- ZrH$_2$. В
$\delta$-ZrH$_{1.6}$ занято около 80\,\% пустот, причём беспорядочно; это и делает $\delta$-фазу устойчивой в широком
интервале составов.
\end{zrbox}

\section{Стереографическая проекция}

Ориентации плоскостей и направлений удобно изображать на плоском рисунке. Поместим кристалл в центр сферы и
проведём из центра нормаль к плоскости до пересечения со сферой --- получим \textbf{полюс} плоскости. Соединим
полюс с южным полюсом сферы; точка, где эта линия пересекает экваториальную плоскость, --- \textbf{стереографическая
проекция} полюса (рис.~\ref{fig:l5_stereo}а). Полюсы северной полусферы попадают внутрь круга, нормали в плоскости
экватора --- на его край.

\fig{l5_stereo}{а) Построение стереографической проекции полюса. б) Стандартная проекция (001) кубического
кристалла: полюсы $\{100\}$, $\{110\}$, $\{111\}$. в) Проекция (0001) циркония: базис в центре, полюсы призм
$\{10\bar10\}$ на краю, пирамид $\{10\bar11\}$ --- в 62\textdegree{} от центра, $\{10\bar17\}$ --- в 15\textdegree.}

У стереографической проекции два полезных свойства: она сохраняет углы (окружности на сфере переходят в
окружности на рисунке), и угол между двумя плоскостями можно измерить прямо по проекции с помощью сетки Вульфа.
Набор полюсов всех важных плоскостей при данной ориентации кристалла --- \textbf{стандартная проекция}
(рис.~\ref{fig:l5_stereo}б, в). Если вместо кристалла взять образец, а вместо полюсов одного кристалла ---
полюсы тысяч зёрен, получится полюсная фигура --- основной способ изображать текстуру (глава~\ref{ch:texture}).
Там же будет нужна и равноплощадная проекция (Ламберта), сохраняющая не углы, а площади: в ней густота точек
прямо показывает долю зёрен.

\section{Гексагональная решётка циркония}

Ниже 863\degC{} цирконий имеет \textbf{гексагональную плотноупакованную} (ГПУ) решётку --- $\alpha$-Zr.
Её ячейка --- шестигранная призма (рис.~\ref{fig:f4_hcp}а) с параметрами $a \approx 0.323$~нм в основании и
$c\approx 0.515$~нм по высоте; отношение $c/a\approx1.593$ немного меньше идеального $\sqrt{8/3}=1.633$.
Важнейшие плоскости:
\begin{itemize}
\item \textbf{базисная} $(0001)$ --- перпендикулярна оси $c$;
\item \textbf{призматические} $\{10\bar10\}$ --- боковые грани призмы, параллельны $c$;
\item \textbf{пирамидальные} $\{10\bar1l\}$ --- наклонные. Плоскость $\{10\bar17\}$, около которой лежат
  гидриды, всего на 15\textdegree{} отклонена от базисной.
\end{itemize}

\fig{f4_hcp}{а) Ячейка ГПУ $\alpha$-Zr: базисная плоскость $(0001)$ (синяя), призматическая $\{10\bar10\}$
(красная), оси $a_1$--$a_3$ и $c$. б) Пластинка $\delta$-гидрида на плоскости, близкой к базисной: в плоскости
пластинки гидрид шире матрицы на $\approx4.6\,\%$, поперёк неё --- на $\approx7.2\,\%$.}

\begin{idea}
ГПУ решётка сильно \emph{анизотропна}: свойства вдоль $c$ и поперёк различаются. Тепловое расширение вдоль
$c$ заметно больше, чем в базисной плоскости. Скольжение дислокаций легко идёт по призматическим плоскостям и
трудно --- с компонентой вдоль $c$. Поэтому то, как ориентированы оси $c$ зёрен (текстура, глава~\ref{ch:texture}),
определяет и прочность трубы, и то, как в ней лягут гидриды.
\end{idea}

\section{Гидриды циркония}

У циркония несколько гидридов:
\begin{itemize}
\item $\delta$-ZrH$_{1.5\ldots1.7}$ --- основной в оболочках. Решётка гранецентрированная кубическая (типа
  флюорита), $a\approx0.478$~нм. Водород занимает тетраэдрические пустоты.
\item $\gamma$-ZrH --- гранецентрированная тетрагональная; метастабильная, появляется при быстром охлаждении.
\item $\varepsilon$-ZrH$_2$ --- при очень высоком водороде.
\end{itemize}
$\delta$-гидрид выпадает с ориентационным соотношением $(0001)_\alpha\parallel\{111\}_\delta$,
$\langle11\bar20\rangle_\alpha\parallel\langle110\rangle_\delta$: плотноупакованные плоскости обеих решёток
параллельны. При таком сопряжении гидрид шире матрицы на $\eps_a\approx4.6\,\%$ в базисной плоскости и на
$\eps_c\approx7.2\,\%$ вдоль оси $c$; объём больше на $\approx17\,\%$ (Carpenter, 1973).

Откуда эти числа? Расстояние между соседними атомами циркония в базисной плоскости $\alpha$-Zr равно $a = 0.3232$~нм,
а в плоскости $\{111\}$ гидрида --- $a_\delta/\sqrt2 = 0.338$~нм: разница 4.6\,\%. Расстояние между плотными
слоями в $\alpha$-Zr --- $c/2 = 0.2574$~нм, в гидриде --- $a_\delta/\sqrt3 = 0.276$~нм: разница 7.2\,\%.

\begin{zrbox}
Гидрид, видимый на шлифе как одна линия длиной в десятки--сотни микрометров, на самом деле --- пачка тонких
\emph{микропластинок} толщиной 0.1--1~мкм и длиной в единицы микрометров. Они лежат около $\{10\bar17\}$ внутри зёрен
и цепляются друг за друга, переходя из зерна в зерно. Модель работает именно с микропластинками, а «макрогидрид»
--- результат их коллективного расположения.
\end{zrbox}

\begin{exercise}
Найдите угол между нормалью к плоскости $(10\bar17)$ и осью $c$ в ГПУ с $c/a = 1.593$.
\end{exercise}

\begin{exercise}
Докажите, что доля объёма, занятая шарами, в ОЦК равна $\pi\sqrt3/8$, а в ГЦК --- $\pi/(3\sqrt2)$.
\end{exercise}

\begin{exercise}
Найдите межплоскостные расстояния $d_{0002}$, $d_{10\bar10}$ и $d_{10\bar11}$ для $\alpha$-Zr. Какие плоскости разнесены
дальше --- базисные или призматические? При каком $c/a$ было бы наоборот? (Скольжение легче идёт по плоскостям,
далеко отстоящим друг от друга: глава~\ref{ch:plasticity}.)
\end{exercise}

\begin{exercise}
Сколько водорода (в ppm по массе) содержит гидрид ZrH$_{1.66}$? А предельный ZrH$_2$?
\end{exercise}
